% ============================================================================
CIFAR-10N: Figure 3 Reproduction and HAS Direction Validation
\label{sec:extra_cifar10n_direction}
% This section asks whether the DIRECTION of HAS's confusion (which class it
% flags as the runner-up) agrees with genuine human-judged confusability.
% The finding is mostly negative.

We test whether HAS's closest decision boundary agrees with where human annotators actually
mis-annotate. 
\begin{table}[htbp]
\centering
\caption{HAS closest-boundary matrix (\%): $P(\text{closest\_boundary}=\text{col} \mid \text{clean label}=\text{row})$ }
\label{tab:cifar10n_has_boundary}
\renewcommand{\arraystretch}{1.1}
\resizebox{\textwidth}{!}{%
\begin{tabular}{@{} l rrrrrrrrrr @{}}
\toprule
\textbf{True} & \textbf{plane} & \textbf{auto} & \textbf{bird} & \textbf{cat} & \textbf{deer} & \textbf{dog} & \textbf{frog} & \textbf{horse} & \textbf{ship} & \textbf{truck} \\
\midrule
airplane   & 0.0 & 0.7 & 0.2 & 0.0 & 98.1 & 0.5 & 0.0 & 0.0 & 0.1 & 0.4 \\
automobile & 2.9 & 0.0 & 0.5 & 0.0 & 0.0 & 0.0 & 0.0 & 0.1 & 0.1 & 96.3 \\
bird       & 0.2 & 92.9 & 0.0 & 3.5 & 0.1 & 0.1 & 0.4 & 0.3 & 2.4 & 0.0 \\
cat        & 0.0 & 0.0 & 0.2 & 0.2 & 0.1 & 0.2 & 98.9 & 0.1 & 0.0 & 0.2 \\
deer       & 99.0 & 0.0 & 0.1 & 0.3 & 0.1 & 0.1 & 0.1 & 0.2 & 0.1 & 0.0 \\
dog        & 8.2 & 0.1 & 0.0 & 0.2 & 0.1 & 0.1 & 0.2 & 0.2 & 0.1 & 90.8 \\
frog       & 0.0 & 0.0 & 0.1 & 99.6 & 0.0 & 0.0 & 0.1 & 0.1 & 0.0 & 0.0 \\
horse      & 0.0 & 26.0 & 0.2 & 1.3 & 0.1 & 0.2 & 71.9 & 0.0 & 0.1 & 0.1 \\
ship       & 0.1 & 0.1 & 97.2 & 0.7 & 0.0 & 0.1 & 0.0 & 1.7 & 0.0 & 0.1 \\
truck      & 0.3 & 76.7 & 0.0 & 0.1 & 0.0 & 22.6 & 0.0 & 0.0 & 0.0 & 0.2 \\
\bottomrule
\end{tabular}}
\end{table}

In sharp contrast to Figure 3 of the CIFAR-10N paper, HAS's closest-boundary mass is 90--99\%
concentrated on a single target class per row, rather than spread across several plausible
neighbours. Several of these dominant pairs are semantically arbitrary: airplane
$\leftrightarrow$ deer (mutual, 98--99\%), cat $\leftrightarrow$ frog (mutual, 99\%).
automobile$\leftrightarrow$truck (96.3\%/76.7\%) is the one pair that does match genuine human
confusion.


HAS class weight-vector geometry, no input data involved


% explanation: pure analysis of the trained classifier head
% (weights/has_model.pth, F.normalize(has_layer.weight, dim=0), then the
% 10x10 cosine similarity matrix between class columns). Confirms the
% collapse is a property of this trained model's geometry, not of which
% images/labels were evaluated.

\begin{table}[htbp]
\centering
\caption{Cosine similarity between HAS's CIFAR-10 class weight vectors (computed directly from the trained classifier head, no input images involved).}
\label{tab:cifar10n_wcos}
\renewcommand{\arraystretch}{1.1}
\resizebox{\textwidth}{!}{%
\begin{tabular}{@{} l rrrrrrrrrr @{}}
\toprule
 & \textbf{plane} & \textbf{auto} & \textbf{bird} & \textbf{cat} & \textbf{deer} & \textbf{dog} & \textbf{frog} & \textbf{horse} & \textbf{ship} & \textbf{truck} \\
\midrule
airplane   & 1.000 & 0.006 & -0.004 & -0.258 & 0.224 & 0.124 & -0.027 & -0.570 & -0.144 & 0.041 \\
automobile & 0.006 & 1.000 & 0.056 & -0.212 & -0.031 & -0.038 & -0.734 & 0.009 & -0.056 & 0.046 \\
bird       & -0.004 & 0.056 & 1.000 & 0.209 & 0.025 & -0.018 & 0.046 & 0.084 & 0.164 & -0.324 \\
cat        & -0.258 & -0.212 & 0.209 & 1.000 & 0.214 & -0.044 & 0.305 & 0.205 & 0.172 & 0.076 \\
deer       & 0.224 & -0.031 & 0.025 & 0.214 & 1.000 & 0.098 & -0.027 & 0.047 & 0.073 & 0.091 \\
dog        & 0.124 & -0.038 & -0.018 & -0.044 & 0.098 & 1.000 & -0.036 & 0.066 & 0.027 & 0.156 \\
frog       & -0.027 & -0.734 & 0.046 & 0.305 & -0.027 & -0.036 & 1.000 & 0.079 & 0.008 & -0.167 \\
horse      & -0.570 & 0.009 & 0.084 & 0.205 & 0.047 & 0.066 & 0.079 & 1.000 & 0.107 & 0.030 \\
ship       & -0.144 & -0.056 & 0.164 & 0.172 & 0.073 & 0.027 & 0.008 & 0.107 & 1.000 & -0.281 \\
truck      & 0.041 & 0.046 & -0.324 & 0.076 & 0.091 & 0.156 & -0.167 & 0.030 & -0.281 & 1.000 \\
\bottomrule
\end{tabular}}
\end{table}

Notably, ``cat'' is the rank-1 nearest weight-vector neighbour for four different classes (bird,
frog, horse, ship) and rank-2 for deer — i.e.\ cat's weight vector sits in a broad, poorly
separated region of $\mathcal{S}^{63}$, which mechanistically explains why ``cat'' also emerges as
a dominant attractor in two independent, unrelated analyses: the synthetic \texttt{test\_drifted}
corruption matrix (Section~\ref{sec:extra_cifar10_sanity}) and the hardest human-disagreement tail
below. An evenly-separated 10-class hyperspherical arrangement would show every off-diagonal
cosine near a small negative value ($\approx -0.11$); the observed range here ($-0.734$ to
$0.305$) is far more uneven than that.

\begin{table}[htbp]
\centering
\caption{Dominant HAS closest-boundary target on the hardest human-disagreement tail (errors $\geq 2/3$, $n=347$ qualifying instances).}
\label{tab:extra_semantic_direction}
\begin{tabular}{@{} l l @{}}
\toprule
\textbf{True class} & \textbf{Dominant HAS target (\% of tail)} \\
\midrule
airplane   & cat (77.3\%) \\
automobile & cat (85.0\%) \\
bird       & cat (85.5\%) \\
cat        & frog (40.0\%) \\
deer       & cat (87.5\%) \\
dog        & cat (93.3\%) \\
frog       & cat (87.5\%) \\
horse      & cat (64.3\%) \\
ship       & cat (76.5\%) \\
truck      & cat (77.8\%) \\
\bottomrule
\end{tabular}
\end{table}

This independently reproduces the same ``cat attractor'' seen in
Table~\ref{tab:cifar10n_has_boundary} .

% ============================================================================
CIFAR-10 Sanity-Check Pipeline (clean-train / worst-label-eval design)

\label{sec:extra_cifar10_sanity}
% This experiment trains on real clean CIFAR-10 labels, then deliberately
% RE-evaluates the identical 50,000 training images with CIFAR-10N's WORST
% human label attached, to test whether HAS notices the mismatch. This is
% not a real generalization/OOD test (the input images never change). Note
% that "accuracy" numbers here are against a noisy label, not real ground
% truth (real-clean-label accuracy is ~99.8-99.98% for both models, i.e.
% near-perfect, since training was genuinely clean).


\begin{table}[htbp]
\centering
\caption{Per-image drift taxonomy on the clean-train/worst-eval CIFAR-10 pass, $n=50{,}000$.}
\label{tab:extra_cifar10_taxonomy}
\begin{tabular}{@{} l rr @{}}
\toprule
\textbf{Category} & \textbf{Baseline} & \textbf{HAS (margin-based)} \\
\midrule
In-Distribution    & 96.2\% & 96.8\% \\
Data Drift         & 2.9\%  & 4.8\%  (Pure Data Drift) \\
Concept Drift      & 0.8\%  & 3.0\%  (Pure Concept Drift) \\
Full Drift (both)  & 0.0\%  & 0.2\%  \\
\bottomrule
\end{tabular}
\end{table}

\begin{table}[htbp]
\centering
\caption{Model ``accuracy'' recomputed two ways on the same 50,000 images.}
\label{tab:extra_cifar10_accuracy_correction}
\begin{tabular}{@{} l cc @{}}
\toprule
\textbf{Model} & \textbf{vs.\ deliberately-noisy Worst label} & \textbf{vs.\ real clean label} \\
\midrule
Baseline & 59.79\% & 99.98\% \\
HAS      & 59.78\% & 99.80\% \\
\bottomrule
\end{tabular}
\end{table}

\begin{table}[htbp]
\centering
\caption{Direction matrix, clean CIFAR-10 test set — sensible, class-pair-specific confusions (HAS).}
\label{tab:extra_direction_clean}
\begin{tabular}{@{} l l @{}}
\toprule
\textbf{Pair} & \textbf{Off-diagonal mass} \\
\midrule
truck $\to$ automobile & 60.0\% \\
horse $\to$ deer       & 45.3\% \\
cat $\to$ dog          & 40.8\% \\
dog $\to$ cat          & 37.2\% \\
ship $\to$ airplane    & 34.7\% \\
\bottomrule
\end{tabular}
\end{table}

\begin{table}[htbp]
\centering
\caption{Direction matrix, synthetically corrupted \texttt{test\_drifted} set — total collapse toward a single class regardless of true label (HAS).}
\label{tab:extra_direction_drifted}
\begin{tabular}{@{} l r @{}}
\toprule
\textbf{True class} & \textbf{\% drifting to ``cat''} \\
\midrule
airplane   & 89.4\% \\
automobile & 85.7\% \\
bird       & 91.2\% \\
deer       & 88.5\% \\
dog        & 96.6\% \\
frog       & 91.4\% \\
horse      & 74.3\% \\
ship       & 76.1\% \\
truck      & 80.0\% \\
\bottomrule
\end{tabular}
\end{table}

The clean-test matrix (Table~\ref{tab:extra_direction_clean}) shows sensible, well-known CIFAR-10
confusions. The synthetically-corrupted set (Table~\ref{tab:extra_direction_drifted}) instead
collapses every class toward ``cat.''


% ============================================================================
Landscape/ARXPHOTO314: Core Drift-Detection Statistics (pre-fine-tuning)
\label{sec:extra_landscape_core}


Table~\ref{tab:extra_pop_tests} reports population-level statistical tests. Each one asks
``are the training-set distribution and the ARXPHOTO314 custom-set distribution the same, or
have they shifted?''

\begin{itemize}
  \item \textbf{MMD (RBF)} — Maximum Mean Discrepancy: compares the two models' 64-D latent
  embedding distributions (train vs.\ custom). Tests data drift directly in feature space,
  independent of labels or predictions.
  \item \textbf{KS (confidence)} — Kolmogorov--Smirnov test on the max-softmax confidence
  distribution (train vs.\ custom). Tests whether output certainty has shifted — a
  concept-drift proxy.
  \item \textbf{KS (HAS margin)} — same KS idea, but on HAS's angular margin score instead of
  confidence.
  \item \textbf{$\chi^2$ (boundary dir.)} — chi-squared test on which class each embedding's
  closest boundary is, comparing the custom set's pattern against the training set's.
\end{itemize}

Baseline only appears in the first two rows because a plain softmax classifier has no geometric
notion of a ``margin'' or a ``closest decision boundary.'' Those concepts only exist because HAS's
angular-margin design puts class weight vectors on a unit hypersphere with actual geometric
boundaries between them. So HAS gets two tests Baseline structurally cannot have at all — the
concrete, measurable version of the paper's core claim that HAS has a geometric signal the
Baseline doesn't.

On the two tests both models share: Baseline's MMD statistic (0.0203) is larger than HAS's
(0.0156) — Baseline's raw, unconstrained embedding space shows more apparent separation between
train and custom. This is likely partly a scale artifact: HAS's embeddings are forced onto a unit
hypersphere ($\ell_2$-normalized), which compresses the range MMD's RBF kernel is sensitive to,
relative to Baseline's unconstrained Euclidean space — so this comparison probably shouldn't be
read as ``HAS drifts less,'' just ``drifts differently, on a different geometric scale.''
Conversely, HAS's KS-confidence statistic (0.4109) is slightly larger than Baseline's (0.4006):
HAS's confidence collapses somewhat more sharply under drift than Baseline's does, plausibly
because angular-margin training pushes correctly-classified training examples to very high,
sharply-peaked confidence, so when confidence does degrade under drift the drop reads as more
discriminative/statistically pronounced.

\begin{table}[htbp]
\centering
\caption{Population-level drift tests, Landscape train vs. ARXPHOTO314 custom set (pre-fine-tuning). All significant at $p<0.01$.}
\label{tab:extra_pop_tests}
\begin{tabular}{@{} l l r r c @{}}
\toprule
\textbf{Test} & \textbf{Model} & \textbf{Statistic} & \textbf{$p$-value} & \textbf{Drifted} \\
\midrule
MMD (RBF)                & Baseline & 0.0203 & $3.32\times10^{-3}$ & \checkmark \\
MMD (RBF)                & HAS      & 0.0156 & $3.32\times10^{-3}$ & \checkmark \\
KS (confidence)          & Baseline & 0.4006 & $\approx 0$ & \checkmark \\
KS (confidence)          & HAS      & 0.4109 & $\approx 0$ & \checkmark \\
KS (HAS margin)          & HAS      & 0.4061 & $\approx 0$ & \checkmark \\
$\chi^2$ (boundary dir.) & HAS      & 42.93  & $1.1\times10^{-8}$ & \checkmark \\
\bottomrule
\end{tabular}
\end{table}

\paragraph{How the taxonomy is generated.} Every sample gets two independent binary
flags, and the four categories below are just the 2$\times$2 combination of the two:
\begin{itemize}
  \item[] \textbf{data-drifted} — computed purely from the 64-D latent embedding (the
  feature vector produced before the classifier head), as its Euclidean distance from the
  training-set centroid, thresholded at $\mu_{\mathrm{dist}} + k\sigma_{\mathrm{dist}}$ using
  training-set statistics (Equation~\ref{eq:thresh} applied to the distance score). This never
  touches softmax output at all, and is computed identically regardless of whether the
  concept-drift half below uses confidence or margin.
  \item[] \textbf{concept-drifted} — computed from a downstream, per-scheme output signal: low
  softmax confidence $q_i$ (Baseline and HAS-conf) or small angular margin $\Delta_i$
  (HAS-margin), each thresholded the same $k$-sigma way (Equation~\ref{eq:thresh}).
\end{itemize}
\begin{center}
\begin{tabular}{@{} l c c @{}}
\toprule
 & concept-drifted = False & concept-drifted = True \\
\midrule
data-drifted = False & In-Distribution & Concept Drift (pure) \\
data-drifted = True  & Data Drift (pure) & Full Drift (both) \\
\bottomrule
\end{tabular}
\end{center}

Table~\ref{tab:extra_taxonomy} compares three ways of sorting the same 2150 custom-set images
into these four categories. HAS-conf and HAS-margin use the exact same trained HAS model and the
exact same data-drift criterion above (latent distance from the training centroid) — the only
thing that differs between them is which signal is used for the concept-drift half: HAS-conf uses
low softmax confidence (scored the same way Baseline is, just on HAS's own confidence output),
while HAS-margin uses the small/negative angular margin
$\Delta_i = \cos_{\mathrm{true}} - \cos_{\mathrm{2nd}}$, HAS's purpose-built geometric signal.

Why the ``Data Drift'' percentages still differ even though the data-drift test is
identical: the table only shows the mutually-exclusive categories, and ``Data Drift'' means
\textbf{pure} data drift (data-drifted \textbf{and not} concept-drifted). Adding back Full Drift
recovers the true total fraction with data-drifted$=$True, which is the same under both
schemes (up to rounding): $15.6\%+28.4\%=44.0\%$ (HAS-conf) vs.\ $4.8\%+39.3\%=44.1\%$
(HAS-margin). What changes is only how that fixed $\sim\!44\%$ splits between ``pure'' and
``full,'' because that split depends on which of those data-drifted samples are also
concept-drifted, and margin flags more samples as concept-drifted overall (total
concept-drifted: $0.0\%+28.4\%=28.4\%$ for HAS-conf vs.\ $2.0\%+39.3\%=41.3\%$ for HAS-margin),
pulling more of the same fixed data-drifted population out of ``pure'' and into ``full.''

\begin{table}[htbp]
\centering
\caption{Drift taxonomy comparison, percentage of custom-set samples per category, three detection schemes. 
The two summary rows make the invariant above explicit: total data-drifted is the same ($\approx$44\%) for 
HAS-conf and HAS-margin: only the pure/full split differs.}
\label{tab:extra_taxonomy}
\begin{tabular}{@{} l r r r @{}}
\toprule
\textbf{Category} & \textbf{Baseline} & \textbf{HAS-conf} & \textbf{HAS-margin} \\
\midrule
In-Distribution      & 72.0\% & 56.0\% & 54.0\% \\
Data Drift (pure)    &  1.8\% & 15.6\% &  4.8\% \\
Concept Drift (pure) & 26.2\% &  0.0\% &  2.0\% \\
Full Drift (both)    &  0.0\% & 28.4\% & 39.3\% \\
\midrule
\emph{Total data-drifted} (Data Drift + Full Drift)       & \emph{1.8\%} & \emph{44.0\%} & \emph{44.1\%} \\
\emph{Total concept-drifted} (Concept Drift + Full Drift) & \emph{26.2\%} & \emph{28.4\%} & \emph{41.3\%} \\
\bottomrule
\end{tabular}
\end{table}

\begin{table}[htbp]
\centering
\caption{Concept-drift direction matrix (HAS, pre-fine-tuning, custom set).}
\label{tab:extra_direction_landscape}
\begin{tabular}{@{} l rrrrr @{}}
\toprule
\textbf{True class} & \textbf{Coast} & \textbf{Desert} & \textbf{Forest} & \textbf{Glacier} & \textbf{Mountain} \\
\midrule
Coast    & 0.099 & 0.314 & 0.124 & 0.158 & 0.305 \\
Desert   & 0.250 & 0.124 & 0.224 & 0.231 & 0.171 \\
Forest   & 0.135 & 0.190 & 0.128 & 0.112 & \textbf{0.435} \\
Glacier  & 0.168 & 0.168 & 0.170 & 0.222 & 0.272 \\
Mountain & 0.215 & 0.188 & 0.205 & 0.210 & 0.183 \\
\bottomrule
\end{tabular}
\end{table}

Forest $\to$ Mountain is the dominant direction. We can visually see this from the flagged images
(Figure~\ref{fig:extra_concept_drift_examples} below).

\textbf{Post-fine-tuning}, the same matrix looks structurally different — Coast becomes the new
attractor instead of Mountain:

\begin{table}[htbp]
\centering
\caption{Concept-drift direction matrix (HAS, POST-fine-tuning, custom set).}
\label{tab:extra_direction_landscape_postft}
\begin{tabular}{@{} l rrrrr @{}}
\toprule
\textbf{True class} & \textbf{Coast} & \textbf{Desert} & \textbf{Forest} & \textbf{Glacier} & \textbf{Mountain} \\
\midrule
Coast    & 0.213 & 0.213 & 0.032 & 0.158 & 0.383 \\
Desert   & \textbf{0.676} & 0.114 & 0.040 & 0.029 & 0.140 \\
Forest   & 0.220 & 0.038 & 0.152 & 0.022 & 0.568 \\
Glacier  & 0.332 & 0.012 & 0.045 & 0.138 & 0.472 \\
Mountain & 0.398 & 0.059 & 0.084 & 0.111 & 0.348 \\
\bottomrule
\end{tabular}
\end{table}
% This suggests fine-tuning shifts WHICH class absorbs ambiguous predictions
% rather than eliminating the attractor effect.

Visualizing drift direction on real images.

Tables~\ref{tab:extra_direction_landscape}
and~\ref{tab:extra_direction_landscape_postft} are aggregate percentages.
Figure~\ref{fig:extra_concept_drift_examples} grounds the same phenomenon in six actual
custom-set photos.

\begin{figure}[htbp]
    \centering
    \resizebox{\textwidth}{!}{\input{figures/concept_drift_examples.pgf}}
    \caption{The 6 most severe (lowest-margin) concept-drifted custom-set images, cross-class
    only. Each panel is captioned true class $\rightarrow$ closest-boundary class and the
    angular margin $\Delta$. Several destinations look semantically arbitrary rather than
    visually motivated.}
    \label{fig:extra_concept_drift_examples}
\end{figure}


% ============================================================================
Landscape: Fine-Tuning Ablation 
\label{sec:extra_landscape_ablation}


In ablation~1, both models are ranked by the same kind of score
(\texttt{data\_drift\_score}, a latent-distance measure) — but each model ranks its own
pool, restricted to each model's own Pure-Data-Drift-flagged subset
(only 10--103 images total). ``All'' instead uses the full 2150-image custom set for both models.

\begin{table}[htbp]
\centering
\caption{``drift-ranked''/``random'' pools are each model's own Pure-Data-Drift-flagged subset (10--103 images), ``all'' uses the full 2150-image custom set.}
\label{tab:extra_ablation1}
\begin{tabular}{@{} l l r r @{}}
\toprule
\textbf{Model} & \textbf{Mode} & \textbf{ft\_pct} & \textbf{$\Delta$ overall error} \\
\midrule
Baseline & drift-ranked & 25/50/100\% & $+7.58$ / $+7.44$ / $+6.33$ \\
Baseline & random       & 25/50/100\% & $+3.68$ / $+5.40$ / $+5.96$ \\
Baseline & all (full custom set) & 25/50/100\% & $-3.02$ / $-2.97$ / $-3.30$ \\
HAS      & drift-ranked & 25/50/100\% & $0.00$ / $-0.79$ / $-0.74$ \\
HAS      & random       & 25/50/100\% & $-0.37$ / $-0.37$ / $-0.46$ \\
HAS      & all (full custom set) & 25/50/100\% & $-5.58$ / $-5.81$ / $-5.86$ \\
\bottomrule
\end{tabular}
\end{table}

Ablation~2 instead ranks each model by its own native output signal ( confidence for
Baseline, angular margin for HAS ) over the full 2150-image custom set at every budget
(a fairer, better-controlled design than ablation~1):

\begin{table}[htbp]
\centering
\caption{Pool = full 2150-image custom set, ranked by each model's own native score (confidence for Baseline, angular margin for HAS).}
\label{tab:extra_ablation2}
\begin{tabular}{@{} l l r r @{}}
\toprule
\textbf{Model} & \textbf{Mode} & \textbf{ft\_pct} & \textbf{$\Delta$ overall error} \\
\midrule
Baseline & baseline-conf & 25/50/100\% & $-1.25$ / $-1.53$ / $-1.76$ \\
Baseline & random        & 25/50/100\% & $-0.65$ / $-1.25$ / $-1.72$ \\
Baseline & all           & 25/50/100\% & $-0.97$ / $-0.79$ / $-0.88$ \\
HAS      & has-margin    & 25/50/100\% & $-0.65$ / $-0.18$ / $+0.38$ \\
HAS      & random        & 25/50/100\% & $-0.23$ / $-0.32$ / $-0.09$ \\
HAS      & all           & 25/50/100\% & $+0.24$ / $+0.56$ / $+0.14$ \\
\bottomrule
\end{tabular}
\end{table}


% ============================================================================
Landscape: Current Two-Pool Fine-Tuning Study (Pool A = Data Drift, Pool B = Concept Drift)
\label{sec:extra_landscape_pool}


``Source acc.'' below is accuracy on the original Landscape test set, evaluated after
fine-tuning on the custom set i.e.\ a catastrophic-forgetting check.

\begin{table}[htbp]
\centering
\caption{Pool A (Pure Data Drift, 103 images total).}
\label{tab:extra_pool_a}
\begin{tabular}{@{} r r rrr rrr @{}}
\toprule
\textbf{ft\_pct} & \textbf{n} & \multicolumn{3}{c}{\textbf{Custom acc.\ (before/after/$\Delta$)}} & \multicolumn{3}{c}{\textbf{Source acc.\ (before/after/$\Delta$)}} \\
\midrule
0.25 & 26  & 73.02 & 73.53 & $+0.51$ & 87.40 & 86.80 & $-0.60$ \\
0.50 & 52  & 73.02 & 73.58 & $+0.56$ & 87.40 & 87.20 & $-0.20$ \\
1.00 & 103 & 73.02 & 74.19 & $+1.16$ & 87.40 & 86.60 & $-0.80$ \\
\bottomrule
\end{tabular}
\end{table}

\begin{table}[htbp]
\centering
\caption{Pool B (Pure Concept Drift + Full Drift, 886 images total).}
\label{tab:extra_pool_b}
\begin{tabular}{@{} r r rrr rr r @{}}
\toprule
\textbf{ft\_pct} & \textbf{n} & \multicolumn{3}{c}{\textbf{Custom acc.\ (before/after/$\Delta$)}} & \multicolumn{2}{c}{\textbf{Margin (before/after)}} & \textbf{Source $\Delta$} \\
\midrule
0.25 & 222 & 73.02 & 75.44 & $+2.42$ & 0.6055 & 0.5851 & $-0.60$ \\
0.50 & 443 & 73.02 & 76.23 & $+3.21$ & 0.6055 & 0.5594 & $-1.20$ \\
1.00 & 886 & 73.02 & 78.56 & $+5.53$ & 0.6055 & 0.4453 & $-1.80$ \\
\bottomrule
\end{tabular}
\end{table}

Pool B (targeting concept-drifted images) gives 2-5$\times$ the accuracy gain of Pool A at every
matched budget, at the cost of real angular-margin erosion (down to 0.445 at 100\%, from a
training baseline of 0.606) and mild monotonic forgetting on the clean source test set. However,
Pool B also has $8.6\times$ more images available (886 vs.\ 103) — some of its larger gain at
matched percentages reflects a larger absolute pool, not only that concept drift is the
``easier win.''


% ============================================================================
Landscape:Targeted vs.\ Global Fine-Tuning at Fixed Budget
\label{sec:extra_rq1}

\begin{table}[htbp]
\centering
\caption{Targeted (Concept) vs.\ Global (random) vs.\ Targeted (Data) fine-tuning, HAS model, concept-drift held-out accuracy.budgets (50 and 110) are shown.}
\label{tab:extra_rq1}
\begin{tabular}{@{} l r r r @{}}
\toprule
\textbf{Condition} & \textbf{$n$ (actual)} & \textbf{acc.\ before} & \textbf{acc.\ after} \\
\midrule
Global\_All       & 50 & 56.11 & 57.70 \\
Targeted\_Concept & 50 & 56.11 & 58.68 \\
Targeted\_Data    & 50 & 56.11 & 55.01 \\
\midrule
Global\_All       & 110 & 56.11 & 58.31 \\
Targeted\_Concept & 110 & 56.11 & 61.12 \\
Targeted\_Data    & 110 & 56.11 & 57.58 \\
\bottomrule
\end{tabular}
\end{table}

Despite the budget-labeling issue, the qualitative finding is consistent and plausible:
Targeted\_Concept beats Global\_All by roughly 2--3$\times$ the error reduction at matched
budgets, and Targeted\_Data (negative control: fine-tuning on data-drift images to fix
concept-drift error) barely moves concept accuracy at all — the expected result if the two
drift types genuinely need different remediation.
